\documentclass[11pt]{article}
\usepackage{amsmath,amsthm,amssymb,mathtools}
\usepackage[margin=1in]{geometry}
\usepackage{hyperref}
\usepackage{booktabs}
\usepackage{enumitem}

\newtheorem{theorem}{Theorem}[section]
\newtheorem{lemma}[theorem]{Lemma}
\newtheorem{proposition}[theorem]{Proposition}
\newtheorem{corollary}[theorem]{Corollary}
\newtheorem{conjecture}[theorem]{Conjecture}
\theoremstyle{definition}
\newtheorem{definition}[theorem]{Definition}
\newtheorem{remark}[theorem]{Remark}
\newtheorem{example}[theorem]{Example}

% --- Macros (consistent with the four source papers) ---
\newcommand{\Hq}{H_q}
\newcommand{\Sn}{S_n}
\newcommand{\Vlam}{V^{\lambda}}
\newcommand{\Om}{\Omega}
\newcommand{\bla}{\widehat\lambda}
\newcommand{\bmu}{\widehat\mu}
\newcommand{\elh}{\ell(\bla)}
\newcommand{\Rp}{{R'}}
\newcommand{\SYT}{\mathrm{SYT}}
\newcommand{\tr}{\mathrm{tr}}
\newcommand{\Eplus}{E^{+}}
\newcommand{\Eminus}{E^{-}}
\newcommand{\qint}[1]{[#1]_q}
\newcommand{\Trs}{T_{\mathrm{rs}}}
\newcommand{\Tstar}{T^{*}}
\newcommand{\Tunif}{T_{\mathrm{uniform}}}
\newcommand{\supp}{\mathrm{supp}}
\newcommand{\HHom}{\mathrm{HH}}
\newcommand{\SBim}{\mathrm{SBim}}
\newcommand{\Outer}{\mathcal{O}}
\DeclareMathOperator{\im}{im}

\title{A trace-vanishing dichotomy for the Iwahori--Hecke algebra of the symmetric group}
\author{Clio Vega}
\date{17 May 2026}

\begin{document}
\maketitle

\begin{abstract}
For partitions $\lambda = (\bla, 1^q)$ with $\elh \ge 2$ and every row of $\bla$ of length $\ge 2$, we exhibit a specific element $\Om^{(\lambda)}$ of the Iwahori--Hecke algebra $\Hq(\Sn)$ whose trace on the Specht module $\Vlam$ obeys a clean combinatorial dichotomy: $\tr_{\Vlam}(\Om^{(\lambda)}) \equiv 0$ as a rational function of $q$ if and only if a combinatorial threshold $\tau(\bla) := \max(0,\, q + 2\elh - 1 - |\bla|)$ is at least $1$. The proof combines four ingredients: (i) a \emph{multi-row sign-kill meta-theorem} (Pillar~1) placing the image of $\Om^{(\lambda)}$ in $\bigcap_{j} \Eminus_j$; (ii) a \emph{right-absorption identity} $\Om^{(\lambda)}(T_1+1) = (q+1)\Om^{(\lambda)}$; (iii) \emph{per-SYT positivity} of every diagonal entry of $\Om^{(\lambda)}$ in the asymmetric Hoefsmit basis; and (iv) an explicit \emph{two-band SYT construction} $\Tunif$ exhibiting a non-vanishing diagonal entry when $\tau(\bla) = 0$. The dichotomy is verified for all $5957$ shapes with $|\lambda| \le 25$. We close with a conjectural categorical lift in the homotopy category of Soergel bimodules.
\end{abstract}

\tableofcontents

\section{Introduction}\label{sec:intro}

\emph{When does the trace of a non-central element of the Iwahori--Hecke algebra vanish identically as a function of $q$?}  For most algebra elements this question is meaningless or trivial.  For a specific family $\{\Om^{(\lambda)}\}$ indexed by partitions $\lambda$ with at least two rows and a column-$1$ tail, we exhibit a clean combinatorial dichotomy.

The element $\Om^{(\lambda)}$ is built out of products of $S_j := T_j + 1$.  For $\nu \vdash N$ with $\ell = \ell(\nu)$, set $\Rp_k := S_{k-1}S_{k-2}\cdots S_1$ and
\[
   \Om^{(\nu)} \;:=\; \Rp_{\ell+1}\,\Rp_{\ell+2}\cdots \Rp_N.
\]
For $\bla$ with $\elh \ge 2$ rows, all of length $\ge 2$, and $\lambda = (\bla,1^q)$, the threshold
\[
   \tau(\bla) \;:=\; \max\bigl(0,\; q + 2\elh - 1 - |\bla|\bigr)
\]
encodes how much the tail length $q$ exceeds the minimum compatible with the shape of $\bla$.

\begin{theorem}[Main dichotomy]\label{thm:main-dichotomy}
Let $\bla$ be a partition with $\elh \ge 2$ rows, every row of length $\ge 2$, and let $\lambda = (\bla,1^q) \vdash n$ for some $q \ge 1$.  Then, as an identity of rational functions of $q$,
\[
   \tr_{\Vlam}\!\bigl(\Om^{(\lambda)}\bigr) \;\equiv\; 0
   \quad\Longleftrightarrow\quad
   \tau(\bla) \;\ge\; 1.
\]
Equivalently, the trace vanishes identically iff $q \ge |\bla| - 2\elh + 2$.  The dichotomy has been verified on all $5957$ shapes $|\lambda| \le 25$ satisfying the row-length hypothesis.
\end{theorem}

\paragraph{Structure of the paper.}
Section~\ref{sec:prelim} fixes notation: the Iwahori--Hecke algebra $\Hq(\Sn)$, the asymmetric Hoefsmit basis with $b'=1$, and the descent eigenspaces $\Eplus_j, \Eminus_j$.  Section~\ref{sec:omega-tau} defines $\Om^{(\lambda)}$ and $\tau$, and tabulates first examples.  Section~\ref{sec:forward} proves the forward direction ($\tau \ge 1 \Rightarrow \tr = 0$) from Pillar~1 and right-absorption.  Sections~\ref{sec:converse-positivity} and~\ref{sec:converse-existence} prove the converse via per-SYT positivity and an explicit construction $\Tunif$.  Section~\ref{sec:sharpness} proves strong sharpness $B^{(\lambda)} \cap \Eminus_{\tau+1} = \{0\}$, pinning the descent depth at exactly $\tau$.  Section~\ref{sec:per-syt-diagnostics} collects three independent per-SYT diagnostics for $\tau$.  Section~\ref{sec:verify} records the scope of computational verification.  Section~\ref{sec:outlook} states a conjectural categorical lift and the Demazure--Lusztig bridge to $K$-theory of the flag variety.

\paragraph{Related work.}
The seminormal basis of $\Hq(\Sn)$ goes back to Hoefsmit~\cite{hoefsmit1974} and Murphy~\cite{murphy1992seminormal}; the cellular structure is due to Dipper--James~\cite{dipperjames1986blocks}.  Ram~\cite{ram1991frobenius,ram2003Skew} developed Markov traces and Jucys--Murphy product machinery in this setting.  Tolmachov--Zhylinskyi~\cite{tolmachov2025markov} compute traces of products of polynomials in Jucys--Murphy elements in closed form; their formula motivates but does not directly settle the dichotomy because $\Om^{(\lambda)}$ is not a central JM product (see Remark~\ref{rem:ram-jm-shadow}).  The categorical home of our dichotomy is the derived horizontal trace $\HHom_*(\SBim_n)$ studied by Gorsky--Hogancamp--Wedrich~\cite{ghw2020derived}; the relevant Specht-isotypic projectors $P_\lambda$ are Hogancamp's categorified Young symmetrizers~\cite{hogancamp2015symmetrizers,abelhogancamp2015}.  Recent work of Trinh--Williams~\cite{trinhwilliams2026} on Hecke trace formulas via partial Springer resolutions is, to our knowledge, the only published reference citing Tolmachov--Zhylinskyi.

Two routes to per-SYT positivity were investigated and \emph{deflated}.  The first, Goertzen--Williamson's variational characterization of the seminormal basis~\cite{goertzenwilliamson2026variational}, gives the right structural picture at $v = 1$ but does not transfer to generic $q$.  The second, Bai--Gu--Guo--Liu's 1423-avoidance criterion for non-negativity in certain pipe-dream rings~\cite{baigugou2026},  matches our diagonals on small shapes but is a different ring; the apparent bridge does not survive scrutiny.  We comment on both in Section~\ref{sec:outlook}.

\section{Preliminaries}\label{sec:prelim}

\subsection{The Iwahori--Hecke algebra}

Let $\Hq(\Sn)$ denote the Iwahori--Hecke algebra of the symmetric group $S_n$ over the field $\mathbb{Q}(q)$, with generators $T_1, \ldots, T_{n-1}$ and relations
\[
   T_i^2 = (q-1)T_i + q,
   \qquad T_iT_{i+1}T_i = T_{i+1}T_iT_{i+1},
   \qquad T_iT_j = T_jT_i \text{ for } |i-j| \ge 2.
\]
The eigenvalues of each $T_j$ on any finite-dimensional module are $q$ and $-1$.  Set
\[
   S_j := T_j + 1,\qquad j = 1, \ldots, n-1.
\]

\begin{lemma}[Square identity]\label{lem:S-square}
$S_j^2 = (q+1) S_j$ in $\Hq(\Sn)$ for every $j$.
\end{lemma}

\begin{proof}
$(T_j+1)^2 = T_j^2 + 2T_j + 1 = ((q-1)T_j + q) + 2T_j + 1 = (q+1)T_j + (q+1) = (q+1)(T_j+1)$.
\end{proof}

Equivalently, $\frac{1}{q+1} S_j$ is an idempotent (the projector onto the $+q$-eigenspace of $T_j$).

\subsection{The Hoefsmit asymmetric seminormal basis}

For $\nu \vdash N$, the Specht module $V^\nu$ has the Hoefsmit seminormal basis $\{v_T : T \in \SYT(\nu)\}$ in which each $T_j$ acts by an explicit two-dimensional block on each pair $\{T, s_j T\}$.  We use the \emph{asymmetric} normalization with $b' = 1$.  Recall the content $c(r, k) := k - r$ of cell $(r,k)$ and the axial distance $d_j(T) := c_{j+1}(T) - c_j(T)$ from $j$ to $j+1$.  The action of $T_j$ on $v_T$ is
\begin{itemize}[topsep=2pt,itemsep=2pt]
\item[(a)] $d_j(T) = 1$ (same row): $T_j v_T = q v_T$.
\item[(b)] $d_j(T) = -1$ (same column): $T_j v_T = -v_T$.
\item[(c)] $|d_j(T)| \ge 2$ (so $s_j T \in \SYT(\nu)$):
\begin{align*}
   T_j v_T &= a_{d}\, v_T + v_{s_j T},\\
   T_j v_{s_j T} &= (a_d a_{-d} + q)\, v_T + a_{-d}\, v_{s_j T},
\end{align*}
where $d = d_j(T)$ and $a_d := (q-1) q^d/(q^d - 1)$.
\end{itemize}

\subsection{Descent eigenspaces}

For each $j \in \{1, \ldots, N-1\}$, the operator $T_j$ has eigenvalues $\{q, -1\}$ on $V^\nu$, giving the orthogonal decomposition
\[
   V^\nu \;=\; \Eplus_j(V^\nu) \;\oplus\; \Eminus_j(V^\nu),
\]
where
\[
   \Eplus_j := \ker(T_j - q)\!\mid_{V^\nu},
   \qquad
   \Eminus_j := \ker(T_j + 1)\!\mid_{V^\nu}.
\]
Equivalently, $S_j$ acts on $\Eplus_j$ by the scalar $q+1$ and on $\Eminus_j$ by $0$.  In the Hoefsmit basis, $T_1$ is diagonal: $T_1 v_T = q v_T$ if $T$ has entries $1, 2$ in the same row (\emph{descent-$0$ at $1$}), and $T_1 v_T = -v_T$ if in the same column (\emph{descent-$1$ at $1$}).

\begin{remark}[Notational overload]\label{rem:q-overload}
The symbol $q$ plays two distinct roles in this paper.  Throughout, $q$ is both the formal Hecke parameter (so that $\Hq$ is defined over $\mathbb{Q}(q)$) and the tail length of $\lambda = (\bla, 1^q)$.  This coincidence reflects the framework in which the four source papers were developed and we keep it; context disambiguates.  The principal results, like Theorem~\ref{thm:main-dichotomy}, are about the trace as an identity of rational functions in $q$ \emph{at the specific integer value $q$ chosen as the tail length}.
\end{remark}

\section{The element $\Om^{(\lambda)}$ and the threshold $\tau$}\label{sec:omega-tau}

\subsection{Definition of $\Om^{(\nu)}$}

For $\nu \vdash N$ with $\ell = \ell(\nu)$, set
\[
   \Rp_k \;:=\; S_{k-1} S_{k-2} \cdots S_1
   \;=\; (T_{k-1}+1)(T_{k-2}+1)\cdots(T_1+1),
   \qquad \Rp_1 := 1,
\]
and define
\[
   \Om^{(\nu)} \;:=\; \Rp_{\ell+1}^{(\nu)}\, \Rp_{\ell+2}^{(\nu)} \cdots \Rp_N^{(\nu)}
   \;\in\; \Hq(\Sn).
\]
We will write $\Om = \Om^{(\lambda)}$ when $\lambda$ is fixed.

\subsection{The threshold $\tau$}

Throughout the remainder of the paper, $\bla$ denotes a partition with $\elh \ge 2$ rows, every row of length $\ge 2$, and we set
\[
   \lambda := (\bla_1, \ldots, \bla_{\elh}, \underbrace{1, \ldots, 1}_{q})
   \;\vdash\; n,\qquad n = q + |\bla|,
\]
with $\ell(\lambda) = \elh + q$.  The threshold is
\[
   \tau(\bla) \;:=\; \max\bigl(0,\; q + 2\elh - 1 - |\bla|\bigr).
\]
Equivalently, writing $|\bla| - 2\elh = \sum_r (\bla_r - 2)$ for the number of cells of $\bla$ outside its first two columns,
\[
   \tau(\bla) \;=\; \max\bigl(0,\; q - 1 - (|\bla| - 2\elh)\bigr).
\]
The condition $\tau(\bla) \ge 1$ is equivalent to $q \ge |\bla| - 2\elh + 2$.

\begin{example}\label{ex:tau-table}
We tabulate small cases.
\begin{center}
\begin{tabular}{l|l|l}
\toprule
$\bla$ & $\tau(\bla)$ & smallest $q$ with $\tau \ge 1$\\
\midrule
$(2,2)$       & $\max(0, q-1)$ & $q = 2$\\
$(3,2)$       & $\max(0, q-2)$ & $q = 3$\\
$(3,3)$       & $\max(0, q-3)$ & $q = 4$\\
$(2,2,2)$     & $\max(0, q-1)$ & $q = 2$\\
$(3,2,2)$     & $\max(0, q-2)$ & $q = 3$\\
\bottomrule
\end{tabular}
\end{center}
The smallest non-trivial dichotomy instance is $\bla = (2,2)$ between $q = 1$ ($\tau = 0$, $\tr \ne 0$) and $q = 2$ ($\tau = 1$, $\tr = 0$).
\end{example}

\subsection{Statement of the main theorem}

We restate Theorem~\ref{thm:main-dichotomy} for reference.

\begin{theorem}[Main dichotomy, restated]\label{thm:dichotomy-restate}
With hypotheses on $\bla$ as above and $\lambda = (\bla, 1^q)$, $q \ge 1$,
\[
   \tr_{\Vlam}(\Om^{(\lambda)}) \equiv 0 \text{ in } \mathbb{Q}(q)
   \iff
   \tau(\bla) \ge 1.
\]
\end{theorem}

The forward direction is proved in Section~\ref{sec:forward}; the converse in Sections~\ref{sec:converse-positivity}--\ref{sec:converse-existence}.

\section{Forward direction: structural nilpotence}\label{sec:forward}

The forward direction follows from two ingredients: a structural containment of the image $\Om \Vlam$ inside $\Eminus_1$ (Pillar~1), and an algebraic identity in the ambient Hecke algebra (right-absorption).  Together they show that, in the descent-at-$1$ block decomposition of $\Vlam$, the matrix of $\Om^{(\lambda)}$ is strictly lower triangular --- in particular nilpotent of order at most~$2$, with vanishing trace.

\subsection{Pillar 1: the multi-row meta-theorem}

\begin{theorem}[Pillar 1, multi-row meta-theorem]\label{thm:pillar1}
Let $\bla$ have $\elh \ge 2$ rows, every row of length $\ge 2$, and let $\lambda = (\bla, 1^q)$ with $q \ge |\bla| - 2\elh + 2$ (equivalently $\tau(\bla) \ge 1$).  Then
\[
   \Om^{(\lambda)} \Vlam \;\subseteq\; \bigcap_{j=1}^{\tau(\bla)} \Eminus_j\!\mid_{\Vlam}.
\]
\end{theorem}

% TODO: full proof from 2026-05-15-multi-row-sign-kill-meta.tex

\begin{proof}[Proof sketch]
Strong induction on $|\bla|$.  The two key ingredients are: (i) a \emph{threshold-matching identity} (multi-row form) showing that every length-preserving contributor $\mu_{k,*}$ obtained from the $E^+_{n-1}$ decomposition preserves $\tau_0(\bla) := q + 2\elh - 1 - |\bla|$ exactly, in both the \emph{regular} ($\bla_k \ge 3$) and \emph{shrinking} ($\bla_k = 2$) cases; (ii) a \emph{support rule}: if every $T$ in the support of a subspace has $p(T) \ge P$ (where $p(T)$ is the column-$1$ identity-prefix length), then the subspace lies in $\bigcap_{j \le P-1} \Eminus_j$.  Vanisher pieces (corner-pair and same-row) inflate $\tau_0$ by $2$ and are killed by the IH plus an outer $\Rp_\ell$ factor whose rightmost $S_1$ annihilates the IH-provided $\Eminus_1$ piece.  Base case $\bla = (2,2)$ at $q \ge 2$ matches the May-14 two-row meta-theorem.  See~\cite[Theorem 1]{multirowMay15} for the complete proof.
\end{proof}

\subsection{Right-absorption}

\begin{lemma}[Right-absorption]\label{lem:right-absorb}
As an element of $\Hq(\Sn)$,
\[
   \Om^{(\lambda)} (T_1 + 1) \;=\; (q+1)\, \Om^{(\lambda)}.
\]
\end{lemma}

\begin{proof}
By Lemma~\ref{lem:S-square}, $(T_1 + 1)^2 = (q+1)(T_1+1)$.  The rightmost factor of $\Om^{(\lambda)}$ is $\Rp_n = (T_{n-1}+1)\cdots(T_1+1)$, ending in $(T_1+1)$.  Therefore
\[
   \Rp_n \cdot (T_1 + 1)
   = (T_{n-1}+1)\cdots(T_2+1) \cdot (T_1+1)^2
   = (q+1)\, \Rp_n.
\]
Multiplying on the left by $\Rp_{\ell+1}\cdots \Rp_{n-1}$ gives the claim.
\end{proof}

\begin{remark}\label{rem:right-absorb-universal}
Lemma~\ref{lem:right-absorb} is an identity in $\Hq(\Sn)$, valid on every $\Hq$-module.  In particular, for any $v \in \Eminus_1$ we have $(T_1+1)v = 0$, hence
\[
   (q+1)\,\Om v \;=\; \Om(T_1+1)v \;=\; 0,
\]
so $\Om v = 0$ since $q+1$ is a unit in $\mathbb{Q}(q)$.  Thus $\Om^{(\lambda)}$ \emph{always} kills $\Eminus_1$ regardless of the value of $\tau(\bla)$.
\end{remark}

\subsection{The trace-vanishing corollary}

\begin{corollary}[Forward direction]\label{cor:forward}
Under the hypotheses of Theorem~\ref{thm:pillar1} (in particular $\tau(\bla) \ge 1$),
\[
   \Om^{(\lambda)\,2} \;=\; 0\quad\text{on } \Vlam,
   \qquad\text{and}\qquad
   \tr_{\Vlam}(\Om^{(\lambda)}) \;=\; 0.
\]
\end{corollary}

\begin{proof}
Order the Hoefsmit basis of $\Vlam$ by descent-at-$1$: place descent-$0$ tableaux first, descent-$1$ tableaux second.  Write the matrix $M(\Om)$ of $\Om^{(\lambda)}$ in this ordering as the block matrix
\[
   M(\Om) \;=\; \begin{pmatrix} A & B \\ X & D \end{pmatrix},
\]
where rows/columns are split as $(\Eplus_1, \Eminus_1)$.  We claim $A = 0$, $B = 0$, $D = 0$.
\begin{itemize}[topsep=2pt,itemsep=2pt]
\item For $T$ descent-$0$ (i.e.\ $v_T \in \Eplus_1$): by Theorem~\ref{thm:pillar1}, $\Om v_T \in \Eminus_1$.  Hence the column of $M(\Om)$ at $T$ has zeros in $\Eplus_1$ rows, so $A = 0$.
\item For $T$ descent-$1$ (i.e.\ $v_T \in \Eminus_1$): by Remark~\ref{rem:right-absorb-universal}, $\Om v_T = 0$.  So columns $B$ and $D$ are both zero.
\end{itemize}
We conclude
\[
   M(\Om) \;=\; \begin{pmatrix} 0 & 0 \\ X & 0 \end{pmatrix},
\]
and immediately $M(\Om)^2 = 0$ and $\tr M(\Om) = 0$.
\end{proof}

This proves the forward direction of Theorem~\ref{thm:main-dichotomy}.  Note that we obtained the stronger conclusion $\Om^2 = 0$, not merely $\tr(\Om) = 0$: the trace-vanishing dichotomy is in fact a \emph{nilpotence} dichotomy.

% TODO: expand structural-nilpotence discussion from 2026-05-16-trace-vanishing-from-pillar1.tex

\section{Converse: per-SYT positivity}\label{sec:converse-positivity}

The converse direction requires showing that $\tr_{\Vlam}(\Om^{(\lambda)}) \not\equiv 0$ whenever $\tau(\bla) = 0$.  We work in the asymmetric Hoefsmit basis with $b' = 1$ and prove that every diagonal matrix entry of $\Om^{(\lambda)}$ lies in the non-negative cone $\mathbb{Q}_{\ge 0}(q) := \{N(q)/D(q) : N, D \in \mathbb{Z}_{\ge 0}[q]\}$.  A sum of elements of this cone vanishes iff each summand does, reducing the converse to exhibiting a single non-vanishing diagonal.

\subsection{Hoefsmit positivity}

\begin{lemma}[Hoefsmit positivity]\label{lem:hoefsmit-positivity}
Every entry of $S_i = T_i + 1$ on $V^\nu$ in the asymmetric Hoefsmit basis ($b' = 1$) lies in $\mathbb{Q}_{\ge 0}(q)$.  Explicitly, with $d = d_i(T)$:
\begin{enumerate}[label=(\roman*),topsep=2pt,itemsep=2pt]
\item Diagonal entries: $[S_i]_{T,T} = q+1$ if $d = 1$; $0$ if $d = -1$; $\qint{d+1}/\qint{d}$ if $|d| \ge 2$.
\item Off-diagonal entries (when $|d| \ge 2$): $[S_i]_{s_iT, T} = 1$ and $[S_i]_{T, s_iT} = q\, \qint{d-1}\,\qint{d+1}/\qint{d}^{2}$.  All other off-diagonals vanish.
\end{enumerate}
\end{lemma}

% TODO: full proof from 2026-05-17-converse-per-syt-positivity.tex

\begin{proof}[Proof sketch]
Direct computation from the Hoefsmit action.  For $|d| \ge 2$, $[S_i]_{T,T} = a_d + 1 = q^{d+1} - 1)/(q^d - 1) = \qint{d+1}/\qint{d}$; for $d \le -2$, rewrite via $\qint{-m} = -q^{-m}\qint{m}$ to obtain $q \qint{k-1}/\qint{k} \in \mathbb{Q}_{\ge 0}(q)$.  The off-diagonal computation likewise yields $q\,\qint{d-1}\qint{d+1}/\qint{d}^2$, manifestly in $\mathbb{Q}_{\ge 0}(q)$.  See~\cite[Lemma~1]{conversePerSYT}.
\end{proof}

\begin{corollary}[Product positivity]\label{cor:product-positivity}
$\mathbb{Q}_{\ge 0}(q)$ is closed under sums and products.  Hence every entry of $\Om^{(\lambda)} = \prod_k \Rp_k$ in the asymmetric Hoefsmit basis is in $\mathbb{Q}_{\ge 0}(q)$.
\end{corollary}

\begin{theorem}[Per-SYT positivity]\label{thm:per-syt}
For every $T \in \SYT(\lambda)$, the diagonal coefficient
\[
   p_T(q) \;:=\; [v_T]\bigl(\Om^{(\lambda)} v_T\bigr)
\]
lies in $\mathbb{Q}_{\ge 0}(q)$.  Consequently,
\[
   \tr_{\Vlam}(\Om^{(\lambda)}) \;=\; \sum_{T \in \SYT(\lambda)} p_T(q) \;\in\; \mathbb{Q}_{\ge 0}(q).
\]
\end{theorem}

\begin{proof}
Immediate from Corollary~\ref{cor:product-positivity}: each $p_T$ is a diagonal entry of $\Om^{(\lambda)}$.
\end{proof}

\subsection{Reduction to one non-vanishing diagonal}

\begin{corollary}[Trace vanishes iff every diagonal vanishes]\label{cor:trace-iff-each}
$\tr_{\Vlam}(\Om^{(\lambda)}) \equiv 0$ as a rational function of $q$ iff $p_T(q) \equiv 0$ for every $T \in \SYT(\lambda)$.
\end{corollary}

\begin{proof}
The cone $\mathbb{Q}_{\ge 0}(q)$ consists of rational functions taking non-negative values at every real $q_0 > 0$ that is not a pole.  A finite sum of such functions vanishing identically forces each summand to vanish at every common point of positivity, hence everywhere.
\end{proof}

To complete the converse, we need a single $T \in \SYT(\lambda)$ with $p_T(q) \not\equiv 0$ when $\tau(\bla) = 0$.

\subsection{Descent-$1$ kill and the stay-walk lower bound}

\begin{lemma}[Descent-$1$ kill]\label{lem:descent1-kill}
If $T$ is descent-$1$ at position $1$, then $\Om^{(\lambda)} v_T = 0$ and hence $p_T(q) = 0$.
\end{lemma}

\begin{proof}
The rightmost factor of $\Om^{(\lambda)}$ is $S_1 = T_1 + 1$, and $S_1 v_T = 0$ for $T$ descent-$1$ at $1$ by Lemma~\ref{lem:hoefsmit-positivity}(i) with $d = -1$.
\end{proof}

So only descent-$0$-at-$1$ SYTs can contribute non-trivially.  Define an SYT $T$ to be \emph{column-descent-free} if no $j$ has $j$ and $j+1$ in the same column at adjacent rows (Definition~\ref{def:col-descent} below).

\begin{lemma}[Stay-walk lower bound]\label{lem:stay-walk}
If $T \in \SYT(\lambda)$ is descent-$0$ at $1$ and column-descent-free, then
\[
   p_T(q) \;\ge\; \prod_{j} [S_j]_{T,T}^{e_j} \;>\; 0
   \quad\text{in } \mathbb{Q}_{\ge 0}(q),
\]
where $e_j$ counts occurrences of $S_j$ in $\Om^{(\lambda)}$.
\end{lemma}

% TODO: full proof from 2026-05-17-converse-per-syt-positivity.tex

\begin{proof}[Proof sketch]
Expand $\Om^{(\lambda)} v_T$ as a sum over Hoefsmit walks $T = T_0 \to T_1 \to \cdots \to T_N = T$, with each step weighted by a matrix entry of some $S_j$.  By Corollary~\ref{cor:product-positivity} every walk contributes a non-negative summand.  The \emph{stay walk} $T_0 = \cdots = T_N = T$ contributes $\prod_k [S_{j_k}]_{T,T}$.  Under the descent-$0$-at-$1$ and column-descent-free hypotheses, every diagonal factor $[S_{j_k}]_{T,T}$ is one of $\{q+1\} \cup \{\qint{d+1}/\qint{d} : |d| \ge 2\}$ --- all strictly positive in $\mathbb{Q}_{\ge 0}(q)$.  The total $p_T$ is bounded below by this stay-walk contribution.
\end{proof}

Combined with Lemma~\ref{lem:descent1-kill}, the converse direction reduces to:

\begin{proposition}[Reduction]\label{prop:converse-reduces}
If there exists $\Tstar \in \SYT(\lambda)$ with descent-$0$ at $1$ and no column-descents, then $\tr_{\Vlam}(\Om^{(\lambda)}) \not\equiv 0$.
\end{proposition}

\section{Converse: existence of $T^{*}$}\label{sec:converse-existence}

We close the converse by giving an explicit construction $\Tunif$ of a descent-$0$-at-$1$ column-descent-free SYT of shape $\lambda$, valid for every $\tau(\bla) = 0$ shape.  The construction is uniform across the regimes $q = 1$ and $q \ge 2$ (in contrast to the earlier case-split version that used the row-superstandard tableau for $q = 1$ and an interleaved construction for $q \ge 2$).

\begin{definition}\label{def:col-descent}
An SYT $T$ of shape $\lambda$ has a \emph{column descent at $j$} if the entries $j$ and $j+1$ sit in the same column at adjacent rows.  $T$ is \emph{column-descent-free} if no $j$ is a column descent.
\end{definition}

\subsection{The construction $\Tunif$}

We use $0$-indexed coordinates.  Let $L = \elh$, $M = |\bla| - 2L$, $n = 2L + M + q$.  Set
\[
   k \;:=\; \#\{r \in \{0, \ldots, L-1\} : \bla_r \ge 3\}.
\]

\paragraph{Column-$0$ values.}  For $r \in \{0, \ldots, L+q-1\}$, define
\begin{equation}\label{eq:vr-def}
   v_r \;:=\;
   \begin{cases}
      2r+1, & 0 \le r < k,\\
      n - 2(L+q-1-r), & k \le r \le L+q-1.
   \end{cases}
\end{equation}
Place $\Tunif(r, 0) := v_r$.  These form a band of small odd values $1, 3, \ldots, 2k-1$ at the top of column~$0$ (the ``long-row band'') and a band of large odd values $\ldots, n-4, n-2, n$ at the bottom (the ``tail band'').

\paragraph{Remaining cells.}  Let $w_1 < w_2 < \cdots < w_{L+M}$ be the remaining values of $\{1, \ldots, n\} \setminus \{v_r\}$ in increasing order.  Fill the subdiagram $\bla' := (\bla_0 - 1, \ldots, \bla_{L-1} - 1)$ row-by-row: row $r$ of $\bla'$ receives values $w_{a_r + 1}, \ldots, w_{a_r + \bla_r - 1}$ at columns $1, \ldots, \bla_r - 1$ of $\lambda$, where $a_r := \sum_{i < r}(\bla_i - 1)$.

\begin{example}[$\bla = (3,2)$ at $q = 1$]
Here $L = 2$, $M = 1$, $k = 1$, $n = 6$, $\lambda = (3, 2, 1)$.  Column-$0$ values: $v_0 = 1$, $v_1 = 4$, $v_2 = 6$.  Remaining $\{2,3,5\}$ fill $\bla' = (2,1)$ row-superstandard: row $0$ receives $2, 3$; row $1$ receives $5$.  Resulting tableau:
\[
   \Tunif \;=\;
   \begin{array}{|c|c|c|}
     \hline 1 & 2 & 3 \\\hline 4 & 5 \\\cline{1-2} 6 \\\cline{1-1}
   \end{array}
\]
Rows and columns are strictly increasing; entries $1, 2$ are in row $0$ (descent-$0$ at $1$); no $j, j+1$ pair shares a column.
\end{example}

\subsection{Three elementary lemmas}

The verification that $\Tunif$ has the required properties reduces to three elementary inequalities about $v_r$ and $w_i$.

\begin{lemma}[Column-$0$ gaps]\label{lem:column0-gaps}
$v_0 = 1$, $v_{L+q-1} = n$, and $v_{r+1} - v_r \ge 2$ for all $r$.
\end{lemma}

% TODO: full proof from 2026-05-16-existence-T-uniform-proof.tex

\begin{lemma}[Row condition $v_r < w_{a_r+1}$]\label{lem:row-condition}
For every $r \in \{0, \ldots, L-1\}$, the column-$0$ value $v_r$ is strictly less than the first remaining value $w_{a_r + 1}$ in row $r$.
\end{lemma}

% TODO: full proof from 2026-05-16-existence-T-uniform-proof.tex.  Key step: writing $a_r = M + r$ for $r \ge k$, the value $w_{M+r+1}$ lies in Section S3 of the $w$-sequence and computes to $v_r + 2q - 1$, with slack $2q-1 \ge 1$.

\begin{lemma}[No column descents]\label{lem:no-col-desc}
No two consecutive integers in $\Tunif$ share a column at adjacent rows.
\end{lemma}

% TODO: full proof from 2026-05-16-existence-T-uniform-proof.tex.  Key fact: the only place in $w_1 < w_2 < \cdots$ where consecutive integers occur is the ``Section S2'' middle band $[k, k+M-q+1]$, and the row-boundary index $M + r + 1$ exceeds this by exactly $q \ge 1$.

\subsection{The main existence theorem}

\begin{theorem}[Existence of $\Tstar$]\label{thm:existence-Tstar}
For every $\bla$ with $\elh \ge 2$, each row of length $\ge 2$, and every $q \ge 1$ satisfying $\tau(\bla) = 0$, the tableau $\Tunif$ of shape $\lambda = (\bla, 1^q)$ is a column-descent-free SYT with $\Tunif(0,0) = 1$ and $\Tunif(0,1) = 2$.
\end{theorem}

\begin{proof}
By Lemma~\ref{lem:column0-gaps}, column $0$ is strictly increasing.  By Lemma~\ref{lem:row-condition}, each row is strictly increasing at the column-$0$-to-column-$1$ transition; strict increase within $\bla'$ follows from the increasing $w$-sequence and row-major filling.  Column strictness within $\bla'$ likewise.  Hence $\Tunif$ is a valid SYT.  The values $\Tunif(0,0) = v_0 = 1$ and $\Tunif(0,1) = w_1 = 2$ are immediate.  By Lemma~\ref{lem:no-col-desc}, no column descents.
\end{proof}

\begin{corollary}[Converse direction]\label{cor:converse}
If $\tau(\bla) = 0$, then $\tr_{\Vlam}(\Om^{(\lambda)}) \not\equiv 0$.
\end{corollary}

\begin{proof}
Combine Theorem~\ref{thm:existence-Tstar} with Proposition~\ref{prop:converse-reduces}.
\end{proof}

Together with Corollary~\ref{cor:forward}, this closes the proof of Theorem~\ref{thm:main-dichotomy}.

\section{Sharpness at $\tau + 1$}\label{sec:sharpness}

Pillar~1 (Theorem~\ref{thm:pillar1}) places $\Om \Vlam$ inside
$\bigcap_{j=1}^{\tau} \Eminus_j$.  A natural question is whether this depth
is best possible.  We prove that it is, in the strong form $B^{(\lambda)} \cap \Eminus_{\tau+1} = \{0\}$.

\subsection{Statement}

Write $B^{(\lambda)} := \Om^{(\lambda)} \Vlam$ for the chain image.

\begin{theorem}[Strong sharpness]\label{thm:sharpness}
Under the standing hypothesis ($\elh \ge 2$, every row $\ge 2$, $q \ge |\bla| - 2\elh + 2$),
\[
   B^{(\lambda)} \,\cap\, \Eminus_{\tau+1}\!\mid_{\Vlam} \;=\; \{0\},
\]
equivalently, $S_{\tau+1} := T_{\tau+1} + 1$ is \emph{injective} on $B^{(\lambda)}$.
\end{theorem}

\begin{corollary}\label{cor:depth-exact}
$\sup\{j : B^{(\lambda)} \subseteq \Eminus_j\!\mid_{\Vlam}\} = \tau(\bla)$
exactly.
\end{corollary}

Theorem~\ref{thm:sharpness} is proved in~\cite{sharpnessTauPlus1}; combined
with Pillar~1~\cite{multirowMay15}, it gives the exact descent depth in
Corollary~\ref{cor:depth-exact}.

\subsection{The one-fact proof}\label{sec:sharp-proof}

The proof rides on a single structural fact about the outer product
$\mathcal{O} := \prod_{j=\ell}^{n-2}(T_j + 1)$ appearing in the Pillar~1
reduction theorem, together with a geometric commutativity observation.

\begin{lemma}[Outer product injectivity]\label{lem:outer-injective}
$\mathcal{O}$ is injective on $\Eplus_{n-1}|_{\Vlam}$, with image in
$\Eplus_\ell|_{\Vlam}$.
\end{lemma}

The proof is a chain application of the adjacent-injectivity
identity $\Eminus_j \cap \Eplus_{j+1} = 0$: for $v \in \Eplus_{n-1}$,
$(T_{n-2}+1)v \in \Eplus_{n-2}$ and the restriction $(T_{n-2}+1)|_{\Eplus_{n-1}}$
is injective; iterating down to $j = \ell$ gives the claim.

\begin{lemma}[Commutativity]\label{lem:commutativity}
Under the standing hypothesis, $S_{\tau+1}$ commutes with every factor of
$\mathcal{O}$ and with each contributor Phi-embedding $\Phi_{k,*}$.
\end{lemma}

\begin{proof}[Proof sketch]
The key inequality is
\[
   \ell - (\tau + 1) \;=\; |\bla| - \elh \;\ge\; \elh \;\ge\; 2,
\]
using every-row-$\ge\!2$ and $\elh \ge 2$.  Hence every factor index
$j \ge \ell$ of $\mathcal{O}$ satisfies $|j - (\tau+1)| \ge 2$, and Hecke
generators at distance $\ge 2$ commute.  The Phi-embeddings are
$T_i$-equivariant for $1 \le i \le n-3$, and the same inequality places
$\tau+1$ in this range.
\end{proof}

\begin{proof}[Proof of Theorem~\ref{thm:sharpness} (sketch)]
Strong induction on $|\bla|$.  Let $w \in B^{(\lambda)}$ with $S_{\tau+1} w = 0$.
By the Pillar~1 reduction theorem,
$w = \mathcal{O}\sum_k \Phi_{k,*}(u_k)$ uniquely with $u_k \in B^{(\mu_{k,*})}$.
By Lemma~\ref{lem:commutativity}, $S_{\tau+1}$ commutes with $\mathcal{O}$
and with each $\Phi_{k,*}$, so
\[
   0 \;=\; S_{\tau+1} w \;=\; \mathcal{O}\!\sum_k \Phi_{k,*}(S_{\tau+1} u_k).
\]
The inner sum lies in $\Eplus_{n-1}|_{\Vlam}$; by Lemma~\ref{lem:outer-injective}
each $\Phi_{k,*}(S_{\tau+1} u_k) = 0$, and disjoint Phi-supports plus
injectivity of each $\Phi_{k,*}$ give $S_{\tau+1} u_k = 0$.  By the
threshold-matching identity, $\tau(\hat\mu_{k,*}) = \tau(\bla)$ for every
contributor, so the IH at $\mu_{k,*}$ gives $u_k = 0$, hence $w = 0$.
\end{proof}

\subsection{Hook base case}\label{sec:hook-base}

The induction bottoms out on the hook layer $\bla = (k_h)$ (length~$1$
shapes, formally outside the standing hypothesis but reached as
contributor inner shapes).  The hook analogue is

\begin{lemma}[Hook sharpness]\label{lem:hook-sharp-survey}
For $\mu = (k_h, 1^{m_h})$ with $k_h \ge 2$ and $m_h \ge k_h$, $S_{m_h - k_h + 2}$
is injective on $B^{(\mu)}$.
\end{lemma}

\begin{proof}[Proof sketch]
Induction on $k_h$ with base case $k_h = 2$ by direct Hoefsmit
computation: $B^{(2, 1^{m_h})} = \Eplus_{m_h+1}|_{V_\mu}$ is a line, spanned
by $w_0 = v_{T^+} + \beta v_{T^-}$ with $T^+$ placing the top two values
at corners $(1,2)$ and $(m_h+1, 1)$ and $T^-$ swapping them.  In $T^+$,
the cells of $m_h$ and $m_h+1$ are non-adjacent (axial distance $m_h \ne \pm 1$),
giving a non-degenerate $T_{m_h}$-block with nonzero off-diagonal; in
$T^-$, same column, so $S_{m_h} v_{T^-} = 0$.  The image $S_{m_h} w_0$
therefore has a nonzero coefficient on a basis vector outside
$\supp(B^{(\mu)})$, witnessing injectivity.  The inductive step uses a
hook analogue of the Pillar~1 reduction with $\Outer^{(\mu)} = \prod_{j=m_h+1}^{m_h+k_h-2}(T_j+1)$
and contributor $\mu_{1,*} = (k_h-1, 1^{m_h-1})$, the same commutativity
argument as above, and the IH on $k_h$.
\end{proof}

\subsection{Verification}\label{sec:sharp-verify}

The full rank $\dim S_{\tau+1}(B^{(\lambda)}) = \dim B^{(\lambda)}$ has been
verified on $14$ datapoints (eight multi-row shapes and six hooks)
computed via a sparse Hoefsmit implementation mod $P = 100003$ at
$q = 11$; see~\cite[Table~1]{sharpnessTauPlus1}.

\section{Per-SYT diagnostics}\label{sec:per-syt-diagnostics}

The converse direction (Section~\ref{sec:converse-positivity}) was proved
by showing each diagonal $p_T(q) := [v_T]\bigl(\Om v_T\bigr)$ lies in the
positive cone $\mathbb{Q}_{\ge 0}(q)$ and exhibiting a non-vanishing
$\Tunif$.  In this section we collect three independent invariants that
pin down which $T$'s contribute and which do not, framing the trace
dichotomy as a confluence of structural diagnostics rather than a single
identity.

\subsection{Diagnostic 1: per-SYT positivity}\label{sec:diag1}

\begin{theorem}[Per-SYT positivity, restated]\label{thm:per-syt-restate}
For every $T \in \SYT(\lambda)$, the diagonal $p_T(q) \in \mathbb{Q}_{\ge 0}(q)$
in lowest terms.
\end{theorem}

This is Theorem~\ref{thm:per-syt}, immediate from Hoefsmit positivity
(Lemma~\ref{lem:hoefsmit-positivity}) and closure of the cone under sums
and products.  Combined with the existence of $\Tunif$
(Theorem~\ref{thm:existence-Tstar}), it closes the converse: $\tau(\bla) = 0$
forces $p_{\Tunif}(q) \ne 0$ on a stay-walk lower bound, so the trace is
non-vanishing.  The cleavage is structural — there is no fortuitous
cancellation between $T$-summands, since none can cancel against another
in a positive cone.

\subsection{Diagnostic 2: leftmost-factor SC vanishing}\label{sec:diag2}

\begin{theorem}[Per-SYT leftmost-factor sign-kill]\label{thm:per-syt-leftmost}
Let $\lambda \vdash n$ with $\ell = \ell(\lambda)$.  If the values $\ell$
and $\ell + 1$ sit in the same column of $T$ (pair $(\ell, \ell+1)$ is SC
at $T$, equivalently $S_\ell v_T = 0$), then $p_T(q) = 0$.
\end{theorem}

\begin{proof}[Two-line proof]
Write $\Om = \Rp_{\ell+1}\Rp_{\ell+2}\cdots \Rp_n = S_\ell \cdot \Rp_\ell \Rp_{\ell+2}\cdots\Rp_n =: S_\ell \cdot X$.
Then $\Om v_T = S_\ell(X v_T) \in \im(S_\ell) = \Eplus_\ell|_{\Vlam}$, while
$v_T \in \Eminus_\ell|_{\Vlam}$ by SC.  Project the seminormal expansion
$\Om v_T = \sum_{T'} c_{T,T'} v_{T'}$ onto $\Eminus_\ell$ along $\Eplus_\ell$:
the left side is $0$, the right side has $c_{T,T''} = 0$ for every
$T'' \in \SYT^{\mathrm{SC}\text{ at }\ell}$; specialise $T'' = T$.
\end{proof}

\begin{remark}
The theorem holds without any restriction on $\lambda$ (multi-row, $\tau$,
staircase).  It is the per-SYT refinement of the universal containment
$B^{(\lambda)} \subseteq \Eplus_\ell|_{\Vlam}$ obtained from the same
leftmost-factor identity.  See~\cite{leftmostPerSYT} for the full
exposition.
\end{remark}

At $\lambda = (3,3,1,1,1)$, an empirical sweep produced $108$ zero diagonals
out of $|\SYT(\lambda)| = 120$.  Theorem~\ref{thm:per-syt-leftmost}
accounts for $26$ of these — exactly the descent-$0$ SYTs with pair
$(5, 6)$ same-column.  In particular, the historically mysterious zero at
\[
   T_0 \;=\; \bigl((1,2,5), (3,4,6), (7), (8), (9)\bigr)
\]
is explained by SC at $(5, 6)$ (value $5$ at row~$0$ column~$2$, value $6$ at
row~$1$ column~$2$).

\subsection{Diagnostic 3: Pillar-1-side SR vanishing}\label{sec:diag3}

The Pillar-1 containment $B^{(\lambda)} \subseteq \Eminus_j$ for
$j = 1, \ldots, \tau$ has a per-SYT shadow obtained by the same projection
argument run at index $j \le \tau$:

\begin{corollary}[Pillar-1-side per-SYT vanishing]\label{cor:pillar1-side-survey}
If $T \in \SYT(\lambda)$ has pair $(j, j+1)$ same-row for some
$j \in \{1, \ldots, \tau\}$ (i.e., $v_T \in \Eplus_j$), then $p_T(q) = 0$.
\end{corollary}

Theorems~\ref{thm:per-syt-leftmost} and Corollary~\ref{cor:pillar1-side-survey}
provide \emph{three} independent invariants pinning $\tau$: structural
depth (Pillar~1 + sharpness), trace vanishing (Theorem~\ref{thm:main-dichotomy}),
and per-SYT leftmost/rightmost diagnostics.  Any future categorical lift
(Conjecture~\ref{conj:categorical}) must unify all three.

\subsection{Open: companion diagnostic and interior zeros}\label{sec:diag-open}

\paragraph{SR companion conjecture.}
Theorem~\ref{thm:per-syt-leftmost} fires at the $\ell$-side via SC.  A
mirror statement is expected at the $\tau$-side: \emph{if pair $(j, j+1)$
is same-row at $T$ for some $j \in \{1, \ldots, \tau\}$, then $p_T(q) = 0$.}
This is exactly Corollary~\ref{cor:pillar1-side-survey} above, so the
``mirror'' is in fact already proved.  A finer companion conjecture
concerns the rightmost-factor identity $\Om = Y \cdot S_1$ extended to
inner positions: a one-page proof is queued in the agent's
proof-attempts file.

\paragraph{Interior D-class zeros.}
The $108 - 26 = 82$ zero diagonals at $\lambda = (3, 3, 1, 1, 1)$ not
explained by Theorem~\ref{thm:per-syt-leftmost} split further: $14$ are
descent-$1$ trivial vanishings (pair $(1, 2)$ SC, $\Om v_T = 0$
outright), and the remaining $68$ are descent-$0$ SYTs with pair
$(\ell, \ell+1)$ diagonal (D-class).  Of these $68$ interior D-class
diagonals, $66$ lie in $\ker \Om$ (column-zero) and are captured by the
two-sided refinement of Diagnostic~2 (see Section~\ref{sec:diag4}
below), so that \emph{only $2$ interior D-class zeros remain} — those
in $\ker \Om^{*} \setminus \ker \Om$ at $\lambda = (3, 3, 1, 1, 1)$.
These two residual zeros are explained by Diagnostic~4 below: they are
\emph{row-zero} vanishings detected by the anti-involution image
$\Om^{*}$, not by $\Om$ itself.  Section~\ref{sec:diag4} closes the
characterisation.

\subsection{Diagnostic 4 --- two-sided kernel dichotomy}\label{sec:diag4}

The third diagnostic captured one half of a Serre-dual pair: it reads
$\Om$ on the left.  A reading on the right --- equivalently, on the
$\iota$-image $\Om^{*}$ under the Hecke anti-involution
$\iota(T_w) = T_{w^{-1}}$ --- closes the characterisation of vanishing
diagonals in the $\tau = 0$ regime.

\begin{theorem}[Two-sided kernel dichotomy]\label{thm:diag4}
Let $\lambda$ with $\tau(\bla) = 0$ and let
$p_T(q) = \langle v_T,\, \Om v_T \rangle$ be the per-SYT diagonal in the
asymmetric Hoefsmit basis.  Then
\[
   p_T(q) \;=\; 0
   \quad\Longleftrightarrow\quad
   \Om v_T = 0 \;\text{ (column-zero)}
   \quad\text{or}\quad
   \Om^{*} v_T = 0 \;\text{ (row-zero)},
\]
where $\Om^{*} := \iota(\Om) = L_n L_{n-1} \cdots L_{\ell+1}$ is the
$\iota$-image of $\Om = \Rp_{\ell+1} \Rp_{\ell+2} \cdots \Rp_n$.
\end{theorem}

\begin{proof}[Proof sketch, forward direction]
The Hecke anti-involution $\iota$ swaps the asymmetric Hoefsmit
parameters $b' = 1$ and $b = 1$ on each generator $\Rp_j$, sending
$\Rp_j = T_j + 1$ to $L_j = T_j + 1$ with the rescaled seminormal
matrix entries (the diagonal axial-distance content is preserved, and
the off-diagonal $b'$-to-$b$ swap is a per-SYT diagonal rescaling).
Consequently $\langle v_T, \Om v_T\rangle = \langle \Om^{*} v_T, v_T\rangle$
in the rescaled pairing, and either $\Om v_T = 0$ (column-zero in the
matrix of $\Om$ at $T$) or $\Om^{*} v_T = 0$ (row-zero) forces
$p_T(q) = 0$.  The full proof is given in
\texttt{2026-05-18-d-class-row-zero.tex}.
\end{proof}

The dichotomy was verified at $\lambda = (3, 3, 1, 1, 1)$ where
$\dim B^{(\lambda)} = 2$ (rank-$2$ compression of $\Om$ forces the
structure): the $78$ D-class zero diagonals split as $76 \in \ker \Om$
(column-zero) and $2 \in \ker \Om^{*} \setminus \ker \Om$ (row-zero), with
no violations across the full $120$ SYTs.  In particular, the two
residual interior D-class zeros of Section~\ref{sec:diag-open} are
exactly the row-zero vanishings: they are invisible to $\Om$ alone but
detected by $\Om^{*}$.

\paragraph{The four diagnostics.}  Combined with the three diagnostics of
Sections~\ref{sec:diag1}--\ref{sec:diag3}, the dichotomy
(Theorem~\ref{thm:main-dichotomy}) is now witnessed by four independent
invariants of $\tau$:

\begin{center}
\begin{tabular}{c l l}
\toprule
\# & Statement & Proof primitive \\
\midrule
1 & $\sup\{j : B^{(\lambda)} \subseteq \Eminus_j\} = \tau$
  & Pillar~1 + adjacent-injectivity \\
2 & $\tr_{\Vlam}(\Om) \equiv 0 \iff \tau \ge 1$
  & $(T_1+1)^2 = (q+1)(T_1+1)$ + Hoefsmit positivity + $\Tunif$ \\
3 & SC/SR pair $\Rightarrow p_T = 0$
  & Leftmost factor + projector identity $\pi_T \circ P^{+}_j = \pi_T$ \\
4 & $p_T = 0 \iff \Om v_T = 0$ or $\Om^{*} v_T = 0$
  & $\iota$-image + $b' = 1 \leftrightarrow b = 1$ rescaling \\
\bottomrule
\end{tabular}
\end{center}

\noindent Diagnostic~4 closes the characterisation of zero diagonals in
the $\tau = 0$ regime up to the verified scope of
Section~\ref{sec:verify}.

\section{Computational verification}\label{sec:verify}

The four ingredients have been verified independently.

\paragraph{Scope.}  The dichotomy was tested on every shape $\lambda = (\bla, 1^q)$ with $|\lambda| \le 25$ satisfying the hypothesis $\elh \ge 2$ and every row $\ge 2$, yielding $5957$ shapes in total: $210$ with $|\lambda| \le 14$; $832$ with $|\lambda| \le 18$; $2688$ with $|\lambda| \le 22$; $5957$ with $|\lambda| \le 25$.

\paragraph{Methodology.}
\begin{itemize}[topsep=2pt,itemsep=2pt]
\item Symbolic computation in \texttt{sympy} verified Hoefsmit positivity (Lemma~\ref{lem:hoefsmit-positivity}) for axial distances $|d| \le 6$.
\item Direct construction confirmed the $\Tunif$ algorithm produces column-descent-free SYTs on all $5957$ shapes ($t\_unified.py$).
\item Pillar~1 and trace-vanishing were verified via sparse application of $\Om^{(\lambda)}$ to random vectors over $\mathbb{F}_p$ with $p = 100003$, $q = 11$, with rank stability checks across multiple probes.
\end{itemize}

\paragraph{Refuted alternatives.}  Three structurally different routes were investigated and ruled out by computational data:
\begin{itemize}[topsep=2pt,itemsep=2pt]
\item A $q$-binomial product ansatz for $\tr(\Om)$ failed at $\lambda = (4,3,1^2)$.
\item The non-crossing matching factorisation (motivated by Jucys--Murphy products) gave rank $12$ where Catalan-counting predicted $5$.
\item The Goertzen--Williamson variational characterisation~\cite{goertzenwilliamson2026variational} of the seminormal basis is only sharp at $v = 1$ (Soergel calculus specialisation), giving no generic-$q$ positivity.
\item The Bai--Gu--Guo--Liu 1423-avoidance criterion~\cite{baigugou2026} matches our diagonals on $7/7$ tested shapes, but the underlying rings differ (theirs is $\mathbb{Z}_{\ge 0}[\beta]$ pipe-dream counts; ours is $\mathbb{Q}_{\ge 0}(q)$ with cyclotomic poles).
\end{itemize}

\begin{remark}\label{rem:ram-jm-shadow}
Ram's content multiset formula~\cite{ram2003Skew} and the Tolmachov--Zhylinskyi closed form for $\tr(\prod p_i(J_i))$~\cite{tolmachov2025markov} suggest a route via Jucys--Murphy products.  However, the multiset $\{c_i(T) : i\}$ is a \emph{shape} invariant, so symmetric JM expressions cannot detect a $\bla$-vs-$q$ dichotomy.  Trace-vanishing of $\Om^{(\lambda)}$ requires a position-asymmetric / non-central JM expression; the right-absorption proof bypasses this entirely.
\end{remark}

\section{Outlook: a categorical lift}\label{sec:outlook}

The forward direction of Theorem~\ref{thm:main-dichotomy} proved strictly more than trace vanishing: it gave $\Om^{(\lambda)\,2} = 0$ on $\Vlam$, with the matrix of $\Om^{(\lambda)}$ in the descent-at-$1$ block decomposition having the form $\bigl(\begin{smallmatrix} 0 & 0 \\ X & 0 \end{smallmatrix}\bigr)$.  This structural rigidity invites a categorification.

\subsection{The conjecture}

\begin{conjecture}[Categorical dichotomy]\label{conj:categorical}
Let $\widetilde\Om^{(\lambda)}$ be a Rouquier complex categorifying $\Om^{(\lambda)}$ as an object of the homotopy category of Soergel bimodules $K^b(\SBim_n)$, and let $P_\lambda$ denote Hogancamp's categorified Young symmetrizer for $\lambda$.  Then
\[
   \tr_{\Vlam}\!\bigl(\Om^{(\lambda)}\bigr) \;\equiv\; 0
   \quad\Longleftrightarrow\quad
   \HHom_*\bigl(P_\lambda \cdot \widetilde\Om^{(\lambda)} \cdot P_\lambda\bigr) \;=\; 0
\]
in bidegree $(q\text{-deg} = |\lambda|,\ \text{hom-deg} = 0)$.
\end{conjecture}

The Specht-isotypic projector $P_\lambda$ realises $\Vlam$ as the $\lambda$-component of the universal Markov trace; the conjecture says that the $\bla$/$\tau$ dichotomy survives intact at the categorical level.

\subsection{What the lift should achieve}

\paragraph{Specht-isotypic decomposition.}  The honest gap is that the Specht-isotypic decomposition
\[
   \mathrm{Tr}(\SBim_n) \;=\; \bigoplus_\lambda \mathrm{Tr}^\lambda
\]
of the horizontal trace into derived blocks is not yet explicitly written down in the derived setting~\cite{ghw2020derived}.  Hogancamp's categorified Young symmetrizers~\cite{hogancamp2015symmetrizers,abelhogancamp2015} provide $P_\lambda$, and the conjectural decomposition should be $\mathrm{Tr}^\lambda = \mathrm{Tr}(\SBim_n \cdot P_\lambda)$.

\paragraph{Lifting right-absorption.}  The identity $S_1^2 = (q+1) S_1$ should categorify to an explicit homotopy / contractibility statement in the categorified bar complex.  This is the categorical analogue of ``$\Om$ ends in $S_1$, and $S_1^2 = (q+1) S_1$'' --- the load-bearing identity behind Lemma~\ref{lem:right-absorb}.

\subsection{Reading list}

Reading list, in suggested order, for tackling Conjecture~\ref{conj:categorical}:
\begin{enumerate}[label=(\arabic*),topsep=2pt,itemsep=2pt]
\item \cite{idempotentsTraceHecke2025} (\texttt{arXiv:2507.10061}) --- Idempotents, traces, and dimensions in Hecke categories.  Sits at the idempotent / trace junction; may already contain a Specht-isotypic decomposition.
\item Trinh--Williams~\cite{trinhwilliams2026} (\texttt{arXiv:2601.17293}) --- Partial Springer resolutions and Hecke trace formulas.
\item Hogancamp~\cite{hogancamp2015symmetrizers} (\texttt{arXiv:1505.08148}) and Abel--Hogancamp~\cite{abelhogancamp2015} (\texttt{arXiv:1510.05330}) --- Categorified Young symmetrizers.
\item Gorsky--Hogancamp--Wedrich~\cite{ghw2020derived} (\texttt{arXiv:2002.06110}) --- Derived horizontal trace of Soergel bimodules.
\item Wedrich's 2025 survey~\cite{wedrich2025survey} (\texttt{arXiv:2509.08478}) for context on the link-homology / TQFT environment.
\end{enumerate}

\subsection{The Demazure--Lusztig bridge}\label{sec:dl-bridge}

A second next-horizon direction is the decategorified bridge between
the LR / Schubert structure constants and the Markov-trace evaluation,
mediated by the Demazure--Lusztig operator on $K_T(G/B)$.  Five published
papers~\cite{amss2019,brubakerBBG2019,brubaker2024,gwzj2025,cheng2025} use
this operator to evaluate equivariant $K$-theory pairings on the flag
variety in terms of Hecke-algebra characters.  The trace
$\tr_{\Vlam}(\Om^{(\lambda)})$ should equal a localisation residue on the
$\lambda$-fixed point of $K_T(G/B)$ via Atiyah--Bott~/ AMSS, and
$\tau(\bla) \ge 1$ should manifest as residue cancellation in this
localisation expansion.  This is the natural next-horizon for a
\emph{structural} (as opposed to algebraic) understanding of $\tau$: the
threshold appears here as the height of a torus-fixed point inside a
Schubert stratification, not as a combinatorial count.  Working out
exactly which Demazure--Lusztig word evaluates to $\Om^{(\lambda)}$ and
extracting the residue cancellation conjecturally encoded by $\tau \ge 1$
is left for future work.

\paragraph{Serre-paired venues.}  Diagnostic~4
(Theorem~\ref{thm:diag4}) is the first to demand internally-paired
(left/right, Serre-dual) structure on the categorical side: the union
$\ker \Om \cup \ker \Om^{*}$ is precisely the kind of zero locus that
lives most naturally inside a Serre-paired bigraded category.  Ho's
``Relative Serre duality for Hecke categories'' (\texttt{arXiv:2504.12798})
is therefore now a primary candidate venue for the categorical lift,
alongside the SBim / Rouquier $\to \mathrm{Coh}(\mathrm{Hilb}_n)^{2\text{-per}}$
framework of Ho--Li (\texttt{arXiv:2305.01306}).

\subsection{What this paper does not address}

We do not address the non-stable regime $q < |\bla| - 2\elh + 2$ where
the standing hypothesis fails, or refinement of
Theorem~\ref{thm:per-syt} to $p_T(q) \in \mathbb{Z}_{\ge 0}[q]$ (the trace
is empirically a non-negative-coefficient polynomial, but our
Hoefsmit-based proof gives only rational positivity); see Remark~3
of~\cite{conversePerSYT}.  The interior D-class zero diagonals at
$\lambda = (3,3,1,1,1)$
(Section~\ref{sec:diag-open}) are now characterised by Diagnostic~4
(Section~\ref{sec:diag4}) as members of $\ker\Om \cup \ker\Om^{*}$.

\begin{thebibliography}{99}

\bibitem{hoefsmit1974}
P.~Hoefsmit.
\emph{Representations of Hecke algebras of finite groups with $BN$-pairs of classical type.}
% TODO: full citation -- PhD thesis, University of British Columbia, 1974.
PhD thesis, 1974.

\bibitem{murphy1992seminormal}
G.~E.~Murphy.
\emph{On the representation theory of the symmetric groups and associated Hecke algebras.}
% TODO: J. Algebra 152 (1992), 492--513.
J.\ Algebra, 1992.

\bibitem{dipperjames1986blocks}
R.~Dipper and G.~James.
\emph{Representations of Hecke algebras of general linear groups.}
% TODO: Proc. London Math. Soc. (3) 52 (1986), 20--52.
Proc.\ London Math.\ Soc., 1986.

\bibitem{ram1991frobenius}
A.~Ram.
\emph{A Frobenius formula for the characters of the Hecke algebras.}
% TODO: Invent. Math. 106 (1991), 461--488.
Invent.\ Math., 1991.

\bibitem{ram2003Skew}
A.~Ram.
\emph{Skew shape representations are irreducible.}
% TODO: in Combinatorial and Geometric Representation Theory, Contemp. Math. 325 (2003), 161--189.
Contemp.\ Math., 2003.

\bibitem{tolmachov2025markov}
K.~Tolmachov and Y.~Zhylinskyi.
\emph{A closed form for the trace of products of Jucys--Murphy elements in the Hecke algebra.}
\texttt{arXiv:2507.19896}, 2025.
% TODO: full bibliographic detail when published.

\bibitem{ghw2020derived}
E.~Gorsky, M.~Hogancamp, and P.~Wedrich.
\emph{Derived traces of Soergel categories.}
\texttt{arXiv:2002.06110}, 2020.
% TODO: published version.

\bibitem{hogancamp2015symmetrizers}
M.~Hogancamp.
\emph{Categorified Young symmetrizers and stable homology of torus links.}
\texttt{arXiv:1505.08148}, 2015.
% TODO: published version.

\bibitem{abelhogancamp2015}
M.~Abel and M.~Hogancamp.
\emph{Categorified Young symmetrizers and stable homology of torus links II.}
\texttt{arXiv:1510.05330}, 2015.
% TODO: published version.

\bibitem{trinhwilliams2026}
M.~T.~Trinh and N.~Williams.
\emph{Partial Springer resolutions and Hecke algebra trace formulas.}
\texttt{arXiv:2601.17293}, 2026.
% TODO: full bibliographic detail.

\bibitem{idempotentsTraceHecke2025}
% TODO: authors.
\emph{Idempotents, traces, and dimensions in Hecke categories.}
\texttt{arXiv:2507.10061}, 2025.

\bibitem{wedrich2025survey}
P.~Wedrich.
\emph{From link homology to TQFTs.}
\texttt{arXiv:2509.08478}, 2025.

\bibitem{goertzenwilliamson2026variational}
% TODO: authors -- Goertzen, Williamson.
\emph{A variational characterisation of the Kazhdan--Lusztig basis.}
\texttt{arXiv:2604.18894}, 2026.

\bibitem{baigugou2026}
Y.~Bai, F.~Gu, Y.~Guo, and Y.~Liu.
\emph{1423-avoidance and positivity in certain pipe-dream rings.}
\texttt{arXiv:2605.10276}, 2026.

\bibitem{multirowMay15}
C.~Vega.
\emph{Extended sign-kill for general multi-row $\bla$: the meta-theorem at arbitrary $\elh \ge 2$.}
Companion preprint, 2026.

\bibitem{traceForward}
C.~Vega.
\emph{Trace vanishing of $\Om^{(\lambda)}$ on $\Vlam$ from the multi-row meta-theorem.}
Companion preprint, 2026.

\bibitem{conversePerSYT}
C.~Vega.
\emph{Per-SYT positivity of $\Om^{(\lambda)}$ in the asymmetric Hoefsmit basis.}
Companion preprint, 2026.

\bibitem{existenceTuniform}
C.~Vega.
\emph{An explicit column-descent-free SYT construction.}
Companion preprint, 2026.

\bibitem{sharpnessTauPlus1}
C.~Vega.
\emph{Sharpness of the multi-row containment at $\tau+1$: $B^{(\lambda)} \cap E^-_{\tau+1} = 0$.}
Companion preprint, 2026.

\bibitem{leftmostPerSYT}
C.~Vega.
\emph{Per-SYT vanishing at the leftmost factor: $p_T(q) = 0$ whenever pair $(\ell, \ell+1)$ is same-column in $T$.}
Companion preprint, 2026.

\bibitem{amss2019}
P.~Aluffi, L.~C.~Mihalcea, J.~Sch\"urmann, and C.~Su.
\emph{Motivic Chern classes of Schubert cells, Hecke algebras, and applications to Casselman's problem.}
\texttt{arXiv:1902.10101}, 2019.

\bibitem{brubakerBBG2019}
B.~Brubaker, V.~Buciumas, D.~Bump, and H.~P.~A.~Gustafsson.
\emph{Colored vertex models and Iwahori Whittaker functions.}
\texttt{arXiv:1906.04140}, 2019.

\bibitem{brubaker2024}
B.~Brubaker, A.~S.~Dasher, M.~Hu, N.~Jain, Y.~Li, Y.~Lin, M.~Mihaila, V.~Tran, and I.~D.~\"Unel.
\emph{Kirillov's conjecture on Hecke--Grothendieck polynomials.}
\texttt{arXiv:2410.07960}, 2024.

\bibitem{gwzj2025}
A.~Gunna, M.~Wheeler, and P.~Zinn-Justin.
\emph{Structure constants for spin Hall--Littlewood functions.}
\texttt{arXiv:2504.19205}, 2025.

\bibitem{cheng2025}
X.~Cheng, R.~Lance, N.~Nagabandi, and A.~Smirnov.
\emph{Quantum difference equations for Grassmannians.}
\texttt{arXiv:2510.21653}, 2025.

\end{thebibliography}

\end{document}
